\section{遺伝子組換え：生態系と食品の別の安全評価}
\label{theory:sustainability:gmo}
新しい形質は生産・医療等の便益へつながり得るが、環境に出た場合の影響も評価する。
カルタヘナ法は2003年公布、2004年施行で、生物多様性への影響を対象とする。
環境中への拡散を伴う第一種使用等と、拡散防止措置を講ずる第二種使用等を区別する。\cite{saf53_cartagena}
食品として食べた場合の健康影響は別の問いである。
厚労省の旧説明では、食品安全委員会の科学的評価を受けて審査し、所定の規格基準の対象を
審査なく流通させない仕組みを示す。\cite{saf53_gmo_food2004,saf53_gmo_standards2000}

生物多様性の承認から食用の安全を自動的に導かず、食用の審査から環境影響がないとも導かない。
対象技術・生物・用途と各手続の適用範囲を確認する。
令和2年度の当時の所管・制度を保持し、後年の行政所管変更や別技術の扱いを遡及しない。
本稿の全例外・全承認手続の照合は未完で、難問査読に残す。
